% !TEX root = ../proyect.tex

\chapter{Arquitectura del sistema y decisiones de diseño}\label{cap:arquitectura}

Este capítulo describe la arquitectura del sistema LinceUS Assistant, las decisiones de diseño adoptadas durante su desarrollo y los diagramas técnicos que representan la estructura de componentes y el modelo de despliegue. El objetivo es ofrecer una visión completa de cómo se organiza internamente el sistema, cómo se comunican sus partes y qué criterios han guiado las principales elecciones técnicas.

% =====================================================
% SECCIÓN 1: ARQUITECTURA GENERAL
% =====================================================

\section{Arquitectura general}\label{sec:arq-general}

LinceUS Assistant sigue una \textbf{arquitectura por capas} (\textit{layered architecture}) con tres niveles claramente diferenciados: capa de presentación, capa de procesamiento conversacional y capa de datos. Esta separación permite que cada capa evolucione de forma independiente y que los cambios en una no afecten directamente a las demás, siempre que se respeten las interfaces definidas entre ellas.

\subsection{Capa de presentación}\label{subsec:capa-presentacion}

La capa de presentación es la interfaz visible para el usuario final. Consiste en un \textbf{widget de chat web} desarrollado en JavaScript puro (vanilla JS), HTML y CSS, diseñado como un componente autocontenido y embebible en cualquier página institucional de la Universidad de Sevilla.

Las principales características de esta capa son:

\begin{itemize}
    \item \textbf{Independencia tecnológica}: El widget no depende de frameworks como React o Angular; se implementa con JavaScript estándar (ES6+), lo que minimiza las dependencias y facilita su integración en páginas web existentes mediante una única etiqueta \texttt{<script>}.
    \item \textbf{Comunicación REST}: El widget se comunica con el servidor Rasa mediante peticiones HTTP POST al endpoint \texttt{/webhooks/rest/webhook}, enviando mensajes del usuario y recibiendo las respuestas del asistente en formato JSON.
    \item \textbf{Renderizado enriquecido}: Soporta formato Markdown en las respuestas del bot, incluyendo tablas, enlaces y listas, para presentar la información académica de forma clara y estructurada.
    \item \textbf{Diseño institucional}: Emplea la identidad visual de la Universidad de Sevilla (colores corporativos, logotipo) para integrarse visualmente en el entorno universitario.
\end{itemize}

\subsection{Capa de procesamiento conversacional}\label{subsec:capa-procesamiento}

Esta capa constituye el núcleo del sistema y está compuesta por varios servicios que colaboran entre sí:

\begin{itemize}
    \item \textbf{Rasa Server}: Motor principal de comprensión del lenguaje natural (NLU) y gestión de diálogo (Core). Recibe los mensajes del usuario, los procesa a través del pipeline NLU para extraer intenciones y entidades, y decide la siguiente acción a ejecutar mediante las políticas de diálogo configuradas.
    \item \textbf{Action Server (Rasa SDK)}: Servidor independiente que ejecuta las acciones personalizadas escritas en Python. Cuando Rasa Core determina que debe ejecutarse una acción personalizada (por ejemplo, consultar la base de datos), delega la ejecución en este servidor a través de una petición HTTP al endpoint \texttt{/webhook}.
    \item \textbf{Google Gemini API}: Servicio externo utilizado para dos propósitos: (1) generación de consultas SQL a partir de lenguaje natural (\textit{Text-to-SQL}) mediante el modelo \texttt{gemma-3-27b-it}, y (2) generación de embeddings vectoriales mediante el modelo \texttt{gemini-embedding-001} para la búsqueda semántica en el pipeline RAG.
    \item \textbf{Ollama} (opcional): Servicio local de inferencia de modelos LLM que actúa como alternativa a Gemini cuando la API externa no está disponible o se agotan las cuotas de uso. Permite ejecutar modelos como \texttt{llama3.2:3b} en CPU sin depender de servicios en la nube.
\end{itemize}

\subsection{Capa de datos}\label{subsec:capa-datos}

La capa de datos se implementa sobre \textbf{Supabase}, una plataforma que proporciona una instancia gestionada de PostgreSQL con extensiones adicionales:

\begin{itemize}
    \item \textbf{PostgreSQL}: Almacena toda la información académica estructurada: universidades, centros, titulaciones, asignaturas, planes docentes y sus metadatos asociados.
    \item \textbf{pgvector}: Extensión que añade el tipo de dato \texttt{vector} y operadores de similitud (distancia coseno, euclidiana y producto punto). Almacena los embeddings de los fragmentos de documentos vectorizados y permite realizar búsquedas semánticas eficientes mediante índices HNSW.
    \item \textbf{Funciones SQL personalizadas}: Funciones almacenadas en la base de datos que encapsulan la lógica de búsqueda vectorial (\texttt{buscar\_plan\_docente()}, \texttt{buscar\_plan\_docente\_global()}), permitiendo filtrar por asignatura, grupo y curso académico con un umbral de similitud configurable.
\end{itemize}

% =====================================================
% SECCIÓN 2: DIAGRAMA DE COMPONENTES
% =====================================================

\section{Diagrama de componentes}\label{sec:diagrama-componentes}

El diagrama de componentes muestra los módulos principales del sistema, sus responsabilidades y las interfaces a través de las cuales se comunican. La Figura~\ref{fig:componentes} representa esta vista.

\begin{figure}[hbtp]
\centering
\resizebox{\textwidth}{!}{%
\begin{tikzpicture}[
    component/.style={rectangle, draw=black, fill=#1, minimum width=3.2cm, minimum height=1.2cm, text width=3cm, align=center, font=\footnotesize\sffamily, rounded corners=2pt},
    component/.default={blue!8},
    layer/.style={rectangle, draw=#1, fill=#1!5, dashed, rounded corners=4pt, inner sep=10pt},
    layer/.default={gray!60},
    arrow/.style={-{Stealth[length=2.5mm]}, thick, #1},
    arrow/.default={black!70},
    label/.style={font=\small\sffamily\bfseries, #1},
    label/.default={black}
]

% === CAPA DE PRESENTACIÓN ===
\node[layer=green!50!black, minimum width=4.5cm, minimum height=2.8cm] (pres) at (0, 0) {};
\node[label=green!50!black, anchor=north west] at ([xshift=3pt, yshift=-3pt] pres.north west) {Presentación};
\node[component=green!12] (widget) at (0, -0.3) {Chat Widget\\(JS/HTML/CSS)};

% === CAPA DE PROCESAMIENTO ===
\node[layer=blue!50!black, minimum width=14cm, minimum height=5.5cm] (proc) at (0, -5) {};
\node[label=blue!50!black, anchor=north west] at ([xshift=3pt, yshift=-3pt] proc.north west) {Procesamiento Conversacional};

\node[component=blue!12] (nlu) at (-4.5, -3.5) {Rasa NLU\\(Pipeline NLP)};
\node[component=blue!12] (core) at (-4.5, -5.3) {Rasa Core\\(Políticas diálogo)};
\node[component=orange!12] (actions) at (0, -4.4) {Action Server\\(Rasa SDK)};
\node[component=orange!12] (txtsql) at (0, -6.2) {Text-to-SQL\\Engine};
\node[component=yellow!15] (gemini) at (4.5, -3.5) {Gemini API\\(LLM + Embeddings)};
\node[component=yellow!15] (ollama) at (4.5, -5.3) {Ollama\\(LLM local)};
\node[component=cyan!12] (rag) at (4.5, -6.8) {RAG Pipeline\\(búsqueda semántica)};

% === CAPA DE DATOS ===
\node[layer=red!40!black, minimum width=14cm, minimum height=2.5cm] (data) at (0, -9.5) {};
\node[label=red!40!black, anchor=north west] at ([xshift=3pt, yshift=-3pt] data.north west) {Capa de Datos};

\node[component=red!10] (pg) at (-4, -9.5) {PostgreSQL\\(Supabase)};
\node[component=red!10] (pgvec) at (0, -9.5) {pgvector\\(embeddings)};
\node[component=red!10] (docs) at (4, -9.5) {Proyectos\\Docentes (PDF)};

% === FLECHAS ===
% Presentación → Procesamiento
\draw[arrow] (widget.south) -- node[right, font=\scriptsize, text=black!60] {REST API} (nlu.north);

% Dentro de procesamiento
\draw[arrow] (nlu.south) -- node[right, font=\scriptsize, text=black!60] {intents, entities} (core.north);
\draw[arrow] (core.east) -- node[above, font=\scriptsize, text=black!60] {webhook} (actions.west);
\draw[arrow] (actions.south) -- (txtsql.north);
\draw[arrow] (actions.east) -- (gemini.west);
\draw[arrow] (actions.east) -- (ollama.west);
\draw[arrow] (gemini.south) -- (rag.north);

% Procesamiento → Datos
\draw[arrow] (txtsql.south) -- (pg.north);
\draw[arrow] (rag.south) -- (pgvec.north);
\draw[arrow] (docs.north) -- node[right, font=\scriptsize, text=black!60] {vectorización} (rag.south west);

\end{tikzpicture}
}%
\caption{Diagrama de componentes del sistema LinceUS Assistant.}
\label{fig:componentes}
\end{figure}

A continuación se describe la responsabilidad de cada componente y las interfaces entre ellos:

\begin{itemize}
    \item \textbf{Chat Widget $\rightarrow$ Rasa NLU}: El widget envía el texto del usuario mediante una petición HTTP POST al endpoint REST de Rasa. La comunicación es síncrona: el widget espera la respuesta para mostrarla en la interfaz.
    \item \textbf{Rasa NLU $\rightarrow$ Rasa Core}: Tras procesar el mensaje a través del pipeline NLU (tokenización, extracción de características, clasificación de intención y extracción de entidades), el resultado se pasa al módulo Core, que decide la siguiente acción.
    \item \textbf{Rasa Core $\rightarrow$ Action Server}: Cuando la acción seleccionada es personalizada (prefijo \texttt{action\_}), Core envía una petición HTTP al Action Server con el tracker completo de la conversación (historial, slots, última intención y entidades).
    \item \textbf{Action Server $\rightarrow$ Text-to-SQL}: Para consultas sobre asignaturas, el Action Server invoca el motor Text-to-SQL, que construye un prompt con el esquema de la base de datos y la pregunta del usuario, lo envía a Gemini y recibe una consulta SQL validada.
    \item \textbf{Action Server $\rightarrow$ Gemini / Ollama}: El Action Server puede invocar Gemini para generación de SQL, detección de titulación o generación de embeddings. Si Gemini no está disponible, se recurre a Ollama como alternativa local.
    \item \textbf{RAG Pipeline $\rightarrow$ pgvector}: La búsqueda semántica genera un embedding de la pregunta del usuario y lo compara contra los embeddings almacenados en pgvector usando similitud coseno, con un umbral mínimo de 0.5.
    \item \textbf{Text-to-SQL $\rightarrow$ PostgreSQL}: Las consultas SQL generadas se ejecutan directamente contra la base de datos PostgreSQL de Supabase, con filtros de seguridad inyectados automáticamente (restricción por titulación).
\end{itemize}

% =====================================================
% SECCIÓN 3: DIAGRAMA DE DESPLIEGUE
% =====================================================

\section{Diagrama de despliegue}\label{sec:diagrama-despliegue}

El diagrama de despliegue describe la distribución física de los componentes del sistema en nodos de ejecución y las conexiones de red entre ellos. La Figura~\ref{fig:despliegue} muestra esta configuración.

\begin{figure}[hbtp]
\centering
\resizebox{\textwidth}{!}{%
\begin{tikzpicture}[
    node/.style={rectangle, draw=black!70, fill=#1, minimum width=4.5cm, minimum height=1cm, text width=4.2cm, align=center, font=\footnotesize\sffamily, rounded corners=2pt},
    node/.default={white},
    host/.style={rectangle, draw=#1, fill=#1!5, rounded corners=4pt, inner sep=12pt, dashed},
    host/.default={gray!60},
    hlabel/.style={font=\small\sffamily\bfseries, #1},
    hlabel/.default={black},
    conn/.style={-{Stealth[length=2.5mm]}, thick, #1},
    conn/.default={black!60},
    biconn/.style={{Stealth[length=2.5mm]}-{Stealth[length=2.5mm]}, thick, #1},
    biconn/.default={black!60}
]

% === NODO: NAVEGADOR ===
\node[host=green!50!black, minimum width=5cm, minimum height=2.5cm] (browser) at (-6, 0) {};
\node[hlabel=green!50!black, anchor=north west] at ([xshift=3pt, yshift=-3pt] browser.north west) {Navegador web};
\node[node=green!10] (wid) at (-6, -0.4) {chatbot-widget.js\\pagina-principal.html};

% === NODO: MÁQUINA LOCAL ===
\node[host=blue!50!black, minimum width=11.5cm, minimum height=6cm] (local) at (0, -5.5) {};
\node[hlabel=blue!50!black, anchor=north west] at ([xshift=3pt, yshift=-3pt] local.north west) {Máquina local (desarrollo)};

\node[node=blue!10] (rasa) at (-3, -3.8) {Rasa Server\\localhost:5005};
\node[node=blue!10] (actserv) at (-3, -5.5) {Action Server\\localhost:5055};
\node[node=orange!10] (ollamanode) at (3, -3.8) {Ollama (opcional)\\localhost:11434};
\node[node=blue!10] (ragpipe) at (3, -5.5) {RAG Pipeline\\(proceso batch)};
\node[node=gray!10] (model) at (0, -7.2) {Modelo Rasa\\(.tar.gz)};

% === NODO: SUPABASE CLOUD ===
\node[host=red!40!black, minimum width=5.5cm, minimum height=3.5cm] (supa) at (-5.5, -11.5) {};
\node[hlabel=red!40!black, anchor=north west] at ([xshift=3pt, yshift=-3pt] supa.north west) {Supabase Cloud (AWS eu-west-3)};
\node[node=red!10] (pgnode) at (-5.5, -11) {PostgreSQL 15\\+ pgvector 0.5};
\node[node=red!10] (apirest) at (-5.5, -12.5) {API REST\\Pooler :6543};

% === NODO: GOOGLE CLOUD ===
\node[host=yellow!60!black, minimum width=5.5cm, minimum height=3.5cm] (gcloud) at (5, -11.5) {};
\node[hlabel=yellow!60!black, anchor=north west] at ([xshift=3pt, yshift=-3pt] gcloud.north west) {Google Cloud};
\node[node=yellow!10] (gemnode) at (5, -11) {Gemini API\\gemma-3-27b-it};
\node[node=yellow!10] (embnode) at (5, -12.5) {Embedding API\\gemini-embedding-001};

% === CONEXIONES ===
\draw[conn] (wid.south) -- node[right, font=\scriptsize, text=black!50] {HTTP POST} ([yshift=5pt] rasa.north);
\draw[conn] (rasa.south) -- node[right, font=\scriptsize, text=black!50] {webhook} (actserv.north);
\draw[biconn, dashed] (actserv.east) -- node[above, font=\scriptsize, text=black!50] {HTTP} (ollamanode.south west);
\draw[conn] (rasa.south east) -- node[right, font=\scriptsize, text=black!50] {} (model.north);

\draw[conn] (actserv.south) -- node[left, font=\scriptsize, text=black!50] {psycopg2 / SQL} (apirest.north);
\draw[conn] (ragpipe.south) -- node[right, font=\scriptsize, text=black!50] {embeddings} (pgnode.north east);
\draw[conn] (actserv.south east) -- node[right, font=\scriptsize, text=black!50, pos=0.4] {HTTPS} (gemnode.north west);
\draw[conn] (ragpipe.south east) -- node[right, font=\scriptsize, text=black!50] {} (embnode.north);

\end{tikzpicture}
}%
\caption{Diagrama de despliegue del sistema LinceUS Assistant.}
\label{fig:despliegue}
\end{figure}

\subsection{Nodos de ejecución}

El sistema se despliega sobre los siguientes nodos:

\begin{enumerate}
    \item \textbf{Navegador web del usuario}: Ejecuta el widget de chat (JavaScript) y se comunica con el servidor Rasa mediante HTTP. No requiere instalación alguna por parte del usuario.

    \item \textbf{Máquina local de desarrollo}: Alberga los tres servicios principales del backend:
    \begin{itemize}
        \item \textit{Rasa Server} (\texttt{localhost:5005}): Sirve la API REST del chatbot y carga el modelo entrenado (\texttt{.tar.gz}).
        \item \textit{Action Server} (\texttt{localhost:5055}): Ejecuta las acciones personalizadas con un timeout de 300 segundos para acomodar consultas LLM en CPU.
        \item \textit{Ollama} (\texttt{localhost:11434}): Servicio opcional de inferencia local de modelos LLM.
        \item \textit{RAG Pipeline}: Proceso batch que se ejecuta bajo demanda para vectorizar nuevos documentos PDF.
    \end{itemize}

    \item \textbf{Supabase Cloud (AWS eu-west-3)}: Instancia gestionada de PostgreSQL con pgvector, accesible mediante connection pooling en el puerto 6543. Aloja tanto los datos estructurados como los embeddings vectoriales.

    \item \textbf{Google Cloud}: Proporciona acceso a la API de Gemini para generación de texto (Text-to-SQL) y generación de embeddings. La comunicación se realiza mediante HTTPS con autenticación por API key.
\end{enumerate}

\subsection{Protocolos de comunicación}

\begin{table}[hbtp]
\centering
\begin{tabular}{|l|l|l|l|}
\hline
\textbf{Origen} & \textbf{Destino} & \textbf{Protocolo} & \textbf{Puerto} \\ \hline
Navegador & Rasa Server & HTTP POST (REST) & 5005 \\ \hline
Rasa Server & Action Server & HTTP POST (webhook) & 5055 \\ \hline
Action Server & Supabase & TCP (psycopg2) & 6543 \\ \hline
Action Server & Gemini API & HTTPS (REST) & 443 \\ \hline
Action Server & Ollama & HTTP (REST) & 11434 \\ \hline
RAG Pipeline & Supabase & TCP (psycopg2) & 6543 \\ \hline
RAG Pipeline & Gemini Embeddings & HTTPS (REST) & 443 \\ \hline
\end{tabular}
\caption{Protocolos y puertos de comunicación entre nodos del sistema.}
\label{tab:protocolos}
\end{table}

% =====================================================
% SECCIÓN 4: DECISIONES DE DISEÑO
% =====================================================

\section{Decisiones de diseño}\label{sec:decisiones-diseno}

A lo largo del desarrollo del sistema se han tomado decisiones de diseño que han condicionado la arquitectura, el rendimiento y la mantenibilidad del proyecto. En esta sección se documentan las más relevantes, indicando el contexto, las alternativas consideradas y la justificación de cada elección.

\subsection{Text-to-SQL con LLM en lugar de reglas estáticas}\label{subsec:dd-texttosql}

\textbf{Contexto:} El sistema debe responder consultas sobre asignaturas que pueden formularse de formas muy diversas (por nombre, curso, tipología, créditos, duración o combinaciones de estos filtros). Implementar todas las variantes mediante reglas \textit{if-else} o expresiones regulares resultaría en un código excesivamente largo, frágil y difícil de mantener.

\textbf{Decisión:} Se optó por un enfoque \textbf{Text-to-SQL basado en LLM}, donde el modelo Gemini (\texttt{gemma-3-27b-it}) recibe un prompt que incluye el esquema de la tabla \texttt{asignaturas}, las columnas permitidas y la pregunta del usuario, y genera la consulta SQL correspondiente.

\textbf{Justificación:}
\begin{itemize}
    \item \textbf{Flexibilidad}: El modelo interpreta preguntas en lenguaje natural sin necesidad de anticipar todas las formulaciones posibles.
    \item \textbf{Mantenibilidad}: Añadir nuevos atributos o filtros solo requiere actualizar el prompt, no reescribir lógica de parsing.
    \item \textbf{Seguridad}: Se aplican mecanismos de protección como una \textit{whitelist} de columnas permitidas, prohibición de subconsultas e inyección automática del filtro de titulación.
\end{itemize}

\textbf{Temperatura 0.0:} Se configura el modelo con temperatura cero para obtener respuestas deterministas, crítico para la generación de SQL donde la variabilidad es indeseable.

\subsection{Gemini como servicio principal con Ollama como alternativa local}\label{subsec:dd-gemini-ollama}

\textbf{Contexto:} El sistema requiere un modelo de lenguaje grande para Text-to-SQL y generación de embeddings. Depender exclusivamente de una API externa implica riesgos de disponibilidad y costes, mientras que depender solo de ejecución local exige hardware potente.

\textbf{Decisión:} Se adoptó una \textbf{estrategia dual}: Gemini como servicio principal (plan gratuito con 15.000 peticiones/día) y Ollama como alternativa local para situaciones donde la API externa no esté disponible.

\textbf{Justificación:}
\begin{itemize}
    \item El plan gratuito de Gemini es suficiente para el volumen de uso académico del proyecto.
    \item Ollama con \texttt{llama3.2:3b} permite seguir operando sin conexión a internet, aunque con menor calidad en las respuestas.
    \item Esta dualidad alinea el proyecto con los requisitos de privacidad (RGPD), ya que el flujo local no envía datos a terceros.
\end{itemize}

\subsection{Embeddings de 2000 dimensiones con HNSW}\label{subsec:dd-embeddings}

\textbf{Contexto:} El modelo \texttt{gemini-embedding-001} genera embeddings de 3072 dimensiones de forma nativa. Sin embargo, el índice HNSW de pgvector presenta limitaciones de rendimiento y almacenamiento con dimensionalidades muy altas.

\textbf{Decisión:} Se redujo la dimensionalidad de salida a \textbf{2000 dimensiones} mediante el parámetro \texttt{output\_dimensionality} de la API de Gemini.

\textbf{Justificación:}
\begin{itemize}
    \item 2000 dimensiones ofrecen un equilibrio entre la calidad de la representación semántica y la eficiencia del índice HNSW (parámetros $m=16$, $ef\_construction=64$).
    \item Reduce el tamaño de almacenamiento por vector en aproximadamente un 35\%.
    \item Las pruebas empíricas muestran que la pérdida de precisión en la búsqueda semántica es negligible para el corpus de documentos del proyecto (planes docentes de la ETSII).
\end{itemize}

\subsection{Chunking con solapamiento para preservar contexto}\label{subsec:dd-chunking}

\textbf{Contexto:} Los planes docentes en PDF contienen información densa que debe fragmentarse para su vectorización. Un chunking demasiado grande diluye la relevancia; uno demasiado pequeño pierde contexto.

\textbf{Decisión:} Se implementó un chunking de \textbf{800 caracteres con 100 caracteres de solapamiento} y un mínimo de 50 caracteres por fragmento.

\textbf{Justificación:}
\begin{itemize}
    \item 800 caracteres ($\approx$150--200 tokens) es suficiente para capturar un párrafo o sección temática completa.
    \item El solapamiento de 100 caracteres garantiza que la información en los límites de fragmentos no se pierda, preservando la coherencia semántica.
    \item Se aplica detección de secciones por expresiones regulares (``Datos básicos'', ``Objetivos'', ``Contenidos'', ``Evaluación'', ``Bibliografía'') para etiquetar cada fragmento con su sección de origen.
\end{itemize}

\subsection{Búsqueda vectorial con fallback a búsqueda por palabras clave}\label{subsec:dd-fallback}

\textbf{Contexto:} La generación de embeddings para las consultas de búsqueda depende de la API de Gemini, que puede agotar su cuota diaria (100 peticiones por minuto, 1000 por día).

\textbf{Decisión:} Se implementó un sistema de \textbf{doble búsqueda}: búsqueda vectorial como método principal y búsqueda por palabras clave (\texttt{ILIKE}) como fallback automático.

\textbf{Justificación:}
\begin{itemize}
    \item La búsqueda vectorial ofrece mejor precisión semántica (encuentra ``sistema de evaluación'' cuando el usuario pregunta ``¿cómo se califica?'').
    \item El fallback por palabras clave garantiza que el sistema siga operativo aunque la API de embeddings no esté disponible, con resultados razonables para consultas que contienen términos exactos.
    \item El umbral de similitud coseno se fija en 0.5, bajo el cual los resultados se descartan por considerarse irrelevantes.
\end{itemize}

\subsection{Fuzzy matching para resolución de nombres}\label{subsec:dd-fuzzy}

\textbf{Contexto:} Los usuarios introducen nombres de asignaturas con errores ortográficos, sin tildes, con abreviaturas o con variaciones coloquiales (``FP'' por ``Fundamentos de Programación'', ``ADDA'' por ``Análisis y Diseño de Datos y Algoritmos'').

\textbf{Decisión:} Se integró la librería \textbf{RapidFuzz} con un umbral de similitud del 80\% combinado con una tabla de alias y acrónimos mantenida en el código.

\textbf{Justificación:}
\begin{itemize}
    \item El algoritmo de distancia de Levenshtein de RapidFuzz tolera errores tipográficos habituales sin requerir entrenamiento adicional.
    \item La tabla de alias cubre las abreviaturas más comunes utilizadas por los estudiantes de la ETSII.
    \item Si la coincidencia difusa falla, se recurre al LLM (Gemini) como último recurso para intentar identificar la asignatura.
\end{itemize}

\subsection{Slots conversacionales para memoria de contexto}\label{subsec:dd-slots}

\textbf{Contexto:} Los usuarios formulan preguntas de seguimiento que hacen referencia implícita a asignaturas o titulaciones mencionadas previamente (``¿Y cuántos créditos tiene?'', ``¿Es obligatoria?'').

\textbf{Decisión:} Se utilizan \textbf{slots de Rasa} para mantener el estado de la conversación:
\begin{itemize}
    \item \texttt{contexto\_titulacion}: Titulación activa seleccionada por el usuario.
    \item \texttt{ultimo\_codigo\_consultado}: Código de la última asignatura consultada.
    \item \texttt{ultimo\_nombre\_asignatura}: Nombre de la última asignatura mencionada.
    \item \texttt{ultimos\_resultados\_asignaturas}: Lista de resultados de la última consulta (para paginación).
\end{itemize}

\textbf{Justificación:}
\begin{itemize}
    \item Los slots permiten que las acciones personalizadas recuperen el contexto sin que el usuario deba repetir información.
    \item La separación entre código y nombre de asignatura facilita tanto las consultas SQL (por código) como las respuestas al usuario (por nombre legible).
    \item La titulación por defecto es GII-IS (Ingeniería del Software), configurable por el usuario mediante la intención \texttt{cambiar\_contexto\_academico}.
\end{itemize}

\subsection{Pipeline NLU de 9 etapas con doble featurización}\label{subsec:dd-pipeline}

\textbf{Contexto:} El pipeline NLU de Rasa debe ser capaz de clasificar correctamente intenciones en español, tolerando errores ortográficos, abreviaturas y variaciones léxicas propias del lenguaje coloquial estudiantil.

\textbf{Decisión:} Se configuró un pipeline de 9 componentes que combina \textbf{featurización a nivel de palabra} (n-gramas 1--2) y \textbf{a nivel de carácter} (n-gramas 1--4):

\begin{enumerate}
    \item \texttt{WhitespaceTokenizer}: Tokenización por espacios en blanco.
    \item \texttt{RegexFeaturizer}: Patrones regex para códigos de asignatura y formatos especiales.
    \item \texttt{LexicalSyntacticFeaturizer}: Características gramaticales del contexto local.
    \item \texttt{CountVectorsFeaturizer} (palabra): N-gramas de palabras (1--2).
    \item \texttt{CountVectorsFeaturizer} (carácter): N-gramas de caracteres (1--4).
    \item \texttt{DIETClassifier}: Clasificador dual de intenciones y entidades (150 épocas, dropout 0.2).
    \item \texttt{EntitySynonymMapper}: Normalización de sinónimos de entidades.
    \item \texttt{ResponseSelector}: Selección de respuestas tipo FAQ.
    \item \texttt{FallbackClassifier}: Detección de baja confianza (umbral 0.7).
\end{enumerate}

\textbf{Justificación:}
\begin{itemize}
    \item La featurización a nivel de carácter es especialmente eficaz para manejar errores ortográficos (``Fundamnetos'' se reconoce por similitud de caracteres con ``Fundamentos'').
    \item El umbral de fallback de 0.7 equilibra la precisión (evitar respuestas incorrectas) con la cobertura (no rechazar demasiadas consultas legítimas).
    \item El \texttt{DIETClassifier} con 150 épocas y dropout de 0.2 proporciona un modelo robusto sin sobreajuste al conjunto de entrenamiento.
\end{itemize}

\subsection{Hash SHA-256 para evitar reprocesamiento de documentos}\label{subsec:dd-hash}

\textbf{Contexto:} El pipeline de vectorización RAG procesa documentos PDF que pueden ser costosos de re-vectorizar (extracción de texto, chunking, generación de embeddings con API de pago por uso).

\textbf{Decisión:} Se almacena un \textbf{hash SHA-256} de cada PDF procesado en la tabla \texttt{planes\_docentes}. Antes de procesar un documento, se compara el hash actual con el almacenado.

\textbf{Justificación:}
\begin{itemize}
    \item Si el hash coincide, el documento no ha cambiado y se omite el reprocesamiento, ahorrando llamadas a la API de embeddings.
    \item Si el hash difiere, se eliminan los fragmentos anteriores y se reprocessa el documento completo, garantizando la consistencia de los datos.
    \item Este mecanismo es especialmente relevante dado el límite de 1000 embeddings por día del plan gratuito de Gemini.
\end{itemize}

% =====================================================
% SECCIÓN 5: PATRONES ARQUITECTÓNICOS
% =====================================================

\section{Patrones arquitectónicos aplicados}\label{sec:patrones}

El diseño del sistema incorpora varios patrones arquitectónicos reconocidos en la ingeniería del software:

\subsection{Arquitectura por capas (\textit{Layered Architecture})}

Como se ha descrito en la Sección~\ref{sec:arq-general}, el sistema se organiza en tres capas con responsabilidades bien definidas. Cada capa solo se comunica con la inmediatamente inferior, lo que reduce el acoplamiento y facilita el reemplazo de componentes (por ejemplo, cambiar el frontend sin afectar la lógica de negocio).

\subsection{Patrón \textit{Action Server} (microservicio de lógica)}

Rasa delega la ejecución de lógica de negocio en un servidor de acciones independiente. Este patrón es similar al patrón \textit{Backend for Frontend} (BFF), donde un servicio intermedio adapta las respuestas de los servicios de datos al formato esperado por el cliente conversacional. La separación permite escalar el Action Server independientemente del motor Rasa.

\subsection{Patrón \textit{Strategy} para clientes LLM}

Los clientes de Gemini (\texttt{gemini\_client.py}) y Ollama (\texttt{ollama\_client.py}) implementan una interfaz similar para la generación de texto. El Action Server puede alternar entre ambos proveedores en función de la disponibilidad, aplicando el patrón \textit{Strategy} para desacoplar la lógica de negocio del proveedor concreto de LLM.

\subsection{Patrón \textit{Pipeline} para procesamiento RAG}

El pipeline de vectorización sigue el patrón \textit{Pipeline} (o \textit{Pipes and Filters}), donde cada etapa transforma los datos y los pasa a la siguiente: extracción de PDF $\rightarrow$ chunking $\rightarrow$ generación de embeddings $\rightarrow$ inserción en base de datos. Cada etapa es independiente y puede modificarse sin afectar a las demás.

\subsection{Patrón \textit{Repository} para acceso a datos}

El módulo \texttt{db.py} centraliza todas las operaciones contra la base de datos, proporcionando una abstracción que oculta los detalles de conexión (pool de conexiones psycopg2) y ejecución de consultas al resto del sistema. Esto facilita el testing unitario mediante mocks y permite cambiar la implementación de persistencia sin modificar la lógica de negocio.
